
\documentclass[letter, 11pt]{article}

    %%%%%%%%%%%%%%%%%%%%%%%%%%%%%%
    \oddsidemargin  0.0in
    \evensidemargin 0.0in
    \textwidth      6.8in
    \headheight     0.0in
    \topmargin      0.0in
    \textheight=9.0in
    \parindent=0in
    %%%%%%%%%%%%%%%%%%%%%%%%%%%%%%

\include{MathCommands}

\usepackage{amsthm,amsfonts,amssymb,amsmath}
\usepackage{graphicx}   %To include graphics
\usepackage[small,compact]{titlesec}

\usepackage[utf8]{inputenc}

\usepackage[francais]{babel}

\usepackage{enumerate}
\usepackage{fancyhdr}
\pagestyle{fancy}

\renewcommand{\headrulewidth}{1pt}
\fancyhead[C]{Automne 2022}
\fancyhead[L]{STT-2100 Régression}
\fancyhead[R]{Audrey-Anne Vallée}

\newcommand{\var}{\mathbb{V}\text{ar}}
\newcommand{\cov}{\mathbb{C}\text{ov}}
\newcommand{\esp}{\mathbb{E}}

\begin{document}

\vspace*{0.05 in}
\begin{center}
{\Large Devoir 1 \\
Dû le 9 octobre 2022 à 23h59}
\end{center}
\vspace*{0.1 in}


À faire individuellement ou en équipe de deux personnes. \\

Vous devez remettre \underline{trois fichiers} par équipe:
\begin{itemize}
\item Un fichier PDF nommé {\tt EquipeX-123.pdf} avec vos réponses aux questions 1 à 3. Vous pouvez simplement numériser une copie papier écrite à la main. Remplacez {\tt X} par le numéro de votre équipe.
\item Un fichier PDF nommé {\tt EquipeX-4.pdf} avec vos réponses de la question 4. Celui-ci doit être rédigé à l'ordinateur et être d'une longueur maximale de 3 pages. %N'indiquez que le numéro de votre équipe sur ce document; certains rapports seront partagés avec les autres étudiants du cours après la remise du devoir.
\item Un document {\tt EquipeX.R}, ou {\tt EquipeX.sas} avec les codes utilisés pour la question 4.\\
\end{itemize}


\underline{La qualité de la langue sera évaluée}, incluant l'orthographe, la grammaire et la syntaxe. Un maximum de 10 points pourra être retranché pour mauvaise qualité du français. \\

Un travail remis en retard ne sera pas accepté. Les étudiants remettant un devoir en retard recevront la note 0.
%Un travail remis entre 17h00 le 15 octobre et 17h00 le 16 octobre sera pénalisé de 20 points. %Un travail remis entre 17h00 le 8 octobre et 17h00 le 9 octobre sera pénalisé de 30 points.
%Aucun travail ne sera accepté après cette date.

\underline{\hspace{17cm}}\\


\textbf{Question 1 (12 pts)}  \\

Dans la section 3.2.1 des notes de cours, nous étudions la distribution de $\widehat{Y}_0$. Nous avons que
$$   \widehat{Y}_0 = \widehat{\beta}_0 + \widehat{\beta}_1 x_0 = \sum_{i=1}^n c_i Y_i, $$
où 
\begin{align*}
    c_i &= \frac{1}{n} - a_i \overline{x} + a_i x_0, \cr
    a_i &= \frac{x_i-\overline{x}}{S_{xx}}.
\end{align*}

\begin{itemize}
\item[(a)] (8 pts)  Montrez que $\displaystyle \sum_{i=1}^n c_i^2 = \frac{1}{n} + \left[ \frac{(\overline{x}-x_0)^2}{S_{xx}} \right].$ Justifiez votre démarche.
\vspace{0.1 in}
\item[(b)] (4 pts) À l'aide du résultat obtenu en (a), montrez que     
$$\mathrm{Var}\left( \widehat{Y}_0 \right) = \left[ \frac{1}{n} + \frac{(\overline{x} - x_0)^2}{S_{xx}} \right] \sigma^2.$$ Justifiez votre démarche.

%\item[(c)] (10 pts) Et si c'était les variables indépendantes qui avaient été mal mesurées? Comment corrigeriez-vous les estimateurs si votre client vous disait que les valeurs de $y_i$ sont correctes, mais que les valeurs de $x_i$ sont toutes surestimées de une unité? Justifiez votre réponse.
\end{itemize}


\newpage

\textbf{Question 2 (20 pts)} 

Soit le modèle de régression
$$
Y_i = \beta_0\,+\,\beta_1\, x_i^2 \,+\, \varepsilon_i\qquad \text{ pour } i=1,\ldots,n.
$$
On fait l'hypothèse que les erreurs $\varepsilon_i$ sont indépendantes et distribuées selon la loi normale avec espérance~0 et variance $\sigma^2$. On prendra pour acquis les trois informations suivantes.

\begin{enumerate}[$\bullet$]\itemsep 3mm
\item L'estimateur de $\beta_1$ par les moindres carrés est $\hat\beta_1= \dfrac{\sum\limits_{i=1}^n x_i^2\,Y_i - \frac{1}{n}\left(\sum\limits_{j=1}^nx_j^2\right)\sum\limits_{i=1}^n Y_i}{\sum\limits_{i=1}^n x_i^4-\frac{1}{n}\left(\sum\limits_{j=1}^nx_j^2\right)^2}$, noté $\hat\beta_1= \dfrac{S_1}{S_2}$.
\item $\hat\beta_1$ est un estimateur sans biais de $\beta_1$.
\item La variance de $\hat\beta_1$ est égale à $\dfrac{\sigma^2}{S_2}$. %$\dfrac{\sigma^2}{\sum\limits_{i=1}^n x_i^4-\frac{1}{n}\left(\sum\limits_{j=1}^nx_j^2\right)^2}$
\end{enumerate}


\begin{enumerate}[(a)]
%\item[(a)] (3 pts) Donnez la distribution de $Y_i$. Justifiez votre réponse.
\item  (7 pts) Calculez l'estimateur des moindres carrés de $\beta_0$. Vous pouvez donner votre réponse en fonction de $\hat\beta_1$ sans le remplacer par sa formule au long. Justifiez votre démarche.
\item (8 pts) Calculez l'espérance de $\hat\beta_0$. Justifiez votre démarche.
\item (8 pts)  Calculez la variance de $\hat\beta_0$. Vous pouvez utiliser le fait que Cov$(\overline Y,\,\hat\beta_1)=0$. Justifiez votre démarche.
%\item   (7 pts) Donnez un intervalle de confiance de niveau ($1-\alpha$) pour $\beta_1$. Justifiez votre réponse.
\end{enumerate}

\vspace{0.2 in}


\textbf{Question 3 (25 pts)} 

Considérons le modèle de régression linéaire simple usuel avec erreurs normales tel que vu en cours. On a déjà déterminé la distribution de $\hat{\beta}_0$ et celle de $\hat{\beta}_1$. On s'intéresse maintenant à la relation entre ces deux estimateurs, ainsi qu'à l'estimation de $\sigma^2$. 

\begin{itemize}\itemsep4mm
\item[(a)] (5 pts) Calculez $\cov (\hat{\beta}_0, \hat{\beta}_1)$. (Vous pouvez utiliser directement tout résultat montré dans les exercices, ainsi que les propriétés générales de la covariance.) Justifiez votre démarche.
%Vous devriez obtenir un résultat qui dépend de $\sigma^2$, $\overline{x}$ et $S_xx$.
%Nous aurons besoin du fait que $\cov(\bar{y}, \hat{\beta_1})$ = 0. Qu'est-ce que ça veut dire concrètement? Est-ce que cela vous semble intuitif? Essayez de justifier le résultat verbalement, sans le démontrer.
\item[(b)] (8 pts) Montrer que $\displaystyle \hat\sigma^2=\frac{1}{n-2} \sum\limits_{i=1}^n \hat\varepsilon_i^2$ est un estimateur sans biais de la variance autour de la droite. (Vous pouvez utiliser  la distribution des résidus abordée dans le chapitre 4.) Justifiez votre démarche.
%Montrez maintenant rigoureusement que $\cov(\bar{y}, \hat{\beta_1})$ = 0.
\item[(c)] (5 pts) Il existe un lien entre la pente d'une droite de régression et la corrélation entre $x$ et $Y$. En justifiant votre démarche, exprimez $\hat\beta_1$ en fonction du coefficient de corrélation échantillonnal $r$ et des écarts-types échantillonnaux  $s_x$ et $s_y$, où
$$ r = \frac{S_{xy}}{  \sqrt{S_{xx}S_{yy}}}; \quad s_x= \sqrt{\frac{1}{n-1}\sum_{i=1}^n (x_i-\overline{x})^2}; \quad s_y= \sqrt{\frac{1}{n-1}\sum_{i=1}^n (Y_i-\overline{Y})^2}.$$
\item[(d)] (7 pts) On s'intéresse à la prédiction du modèle RLS pour une valeur de $x$ située à 3 écarts-types sous la moyenne, donc en $x_0=\overline x-3\,s_x$. En utilisant le résultat trouvé en c), déterminez si la valeur prédite de $Y$ se situera à 3 écarts-types ($s_y$) sous la moyenne $\overline y$. Dans le cas contraire,  combien d'écarts-types $s_y$ séparent la valeur prédite de la moyenne $\overline y$? Justifiez votre démarche.\\({\it Indice: remplacez d'abord $\hat\beta_0$ dans l'équation de la valeur prédite.})

\end{itemize}




\newpage
\vspace{0.2 in}
\textbf{Question 4 (40 pts)} 

Les tuyaux d'égouts municipaux sont soumis à la corrosion et à des bris d'usure.  On s'intéresse à la taille des dommages à la paroi extérieure des égouts. Il est possible de mesurer la taille d'un dommage sur le terrain (sans creuser ni entrer dans l'égout) à l'aide d'un appareil de mesure à ultrasons. On obtiendrait une meilleure mesure en déterrant le tuyau, mais cela implique des opérations coûteuses. %On s'est soumis à l'exercice pour un certain nombre de tuyaux, en espérant obtenir un modèle stochastique prédictif pour l'avenir.
Un modèle linéaire pourrait permettre de prédire la taille véritable d'un dommage  à partir de la mesure ultrasons et ainsi évaluer si un tuyau est très endommagé sans le déterrer. \\

Chaque équipe recevra un jeu de données de 90 mesures de dommages prises sur des tuyaux différents. %, sélectionné au hasard parmi un plus grand jeu de données.
Pour chaque tuyau, on vous donne la taille (en mm) de la détérioration mesurée sur le terrain, et celle (toujours en mm) mesurée en déterrant le tuyau.  Vous devez rédiger des parties d'un rapport d'analyse de régression en répondant aux questions suivantes. Utilisez $\alpha = 0,01$ pour toutes vos inférences. \\

{\bf Dans le fichier PDF}: identifiez clairement la question associée à votre réponse (par exemple (i) a)). Donnez les réponses aux questions et les tableaux de résultats obtenus avec votre logiciel qui sont nécessaires pour
répondre aux questions. Les figures et tableaux inclus dans vos réponses doivent être présentés de façon professionnelle. Ils doivent avoir des titres clairs et pertinents. Chaque figure présentée doit être discutée dans vos réponses. 

\textbf{Sorties statistiques:}
Vous devez fournir les programmes R ou SAS utilisés pour obtenir vos réponses. Un programme qui ne fonctionne pas ou qui est mal annoté entraînera une pénalité. \\


\begin{enumerate}[(i)]
	\item (9 pts) {\bf Décription du jeu de données et de l'objectif poursuivi.}
\begin{enumerate}[a)]
\item Identifiez et définissez les variables qui nous intéressent.
%\item Mentionner le nombre et la source des observations.
\item Décrivez {\it brièvement} la distribution de chacune des variables (graphiques et statistiques descriptives). Commentez la distribution conjointe des variables à l'aide d'un diagramme de dispersion.  (8 lignes maximum, sans compter les figures).
\item Décrire l'objectif de l'analyse de régression. (4 lignes maximum).
\end{enumerate}
	\item (14 pts) {\bf Choix de modèle utilisé.}
\begin{enumerate}[a)]
	%\item Si on choisit de transformer les variables avant de procéder \`a l'analyse, justifier son choix (raisons théoriques, expérience antérieure avec ce type de données, intuition seulement\,?).
	%\item Au besoin, simplifier le mod\`ele (par exemple, en refaisant la régression sans $\beta_0$).
	\item Vérifiez l'adéquation du mod\`ele RLS et le respect des postulats au moyen de graphiques et tests pertinents. Modifiez le modèle de façon à ce que vous soyez satisfait de la justesse des postulats. (On ne peut pas toujours modifier les variables pour que tous les postulats soient exactement vérifiés;  choisissez alors le meilleur modèle possible). Dans le fichier réponse, formulez clairement le modèle et les postulats du mod\`ele retenu (i.e. écrire le modèle complet en language mathématique). Présentez seulement le modèle final, et non tous les modèles ajustés. Motivez le choix de modèle final (8 lignes maximum). Incluez toutes les étapes de votre démarche (différents modèles, vérification des postulats) dans votre programme R ou SAS. 
\end{enumerate}
	\item (17 pts) {\bf Présentation des résultats.}
\begin{enumerate}
	\item Donnez les estimations des param\`etres (${\hat \beta}_0$, ${\hat \beta}_1$, ${\hat \sigma}^2$).
	\item Fournissez des intervalles de confiance pour les paramètres (${\hat \beta}_0$, ${\hat \beta}_1$) et une interprétation en contexte. (8 lignes maximum pour l'interprétation).
	\item Discutez de la signification statistique des coefficients ${\hat \beta}_0$ et ${\hat \beta}_1$ à l'aide de tests $t$. (6 lignes maximum).
	\item L'analyse de variance permet de faire un test d'hypothèses. Donnez les hypothèses de ce test et discutez de sa signification statistique pour vos données. (4 lignes maximum pour la discussion).
\end{enumerate}
\end{enumerate}




\end{document}




\vspace{0.3in}
\textbf{Question 4 (10 pts)} (Design d'expérience)\\

On s'intéresse à un médicament sensé réduire le nombre de crises chez les individus atteints d'epilepsie. Vous devez aider à planifier une expérience pour étudier l'effet de la dose du médicament sur le nombre de crises. Puisque vous vous intéressez à la relation entre deux variables quantitatives, vous planifiez analyser les données avec un modèle de régression linéaire simple. \\

Soit $x_1, \ldots, x_n$ les valeurs de la dose pour les 20 individus de votre étude.\\ Pour chacun, vous pouvez attribuer n'importe quelle dose entre 2 mg et 21 mg. \\

\begin{itemize}
%\item [(a)] Pour le modèle de régression linéaire simple $Y_i = \beta_0 + \beta_1 x_i + \epsilon_i$, prouvez que la variance de $\beta_0$ est minimale si les $x_i$ sont tels que $\bar{x}$ = 0.
\item [(a)] (5 pts) Quelle est la meilleure fa\c{c}on de choisir les valeurs $x_1,\ldots,x_n$ pour minimiser la variance de votre estimateur $\hat{\beta_1}$? Justifiez votre réponse.
\vspace{0.1 in}
\item [(b)] (5 pts) Pourquoi ne serait-ce pas vraiment une bonne idée de choisir les valeurs de $x_i$ de cette façon, même si cela minimise la variance de votre estimateur?
\end{itemize}


\end{document} 